\section{Approximation and computational resources}\label{sec:resources53}
An action-contrast certificate is constructive only if its representation, integration and uniformity costs are included. A full state-pair table is not required by the preceding theorem. Neither does avoiding that table automatically remove a dimensional or horizon dependence.

\begin{proposition}[Certified compressed pair functions]\label{prop:compressed53}
Let $\psi_t(x,y)$ be any measurable approximation to $e_t(x)-e_t(y)$ on domains containing the pairs reached by the chosen couplings. Suppose verified numbers $\eta_t\ge0$ bound the defects
\begin{align}\label{eq:pairdefect53}
 |\psi_T-(r_T(x)-r_T(y))|&\le\eta_T,\notag\\
 |\psi_t-(r_t(x)-r_t(y))-\beta_t\Lambda_t^\pi\psi_{t+1}|&\le\eta_t.
\end{align}
Then, with $a_T=\eta_T$ and $a_t=\eta_t+\beta_ta_{t+1}$,
\begin{equation}\label{eq:pairerror53}
 |e_t(x)-e_t(y)-\psi_t(x,y)|\le a_t.
\end{equation}
If a computed action-coupling integral of $\psi_{t+1}$ has certified error at most $\tau_t$, its center with radius $a_{t+1}+\tau_t$ encloses the action contrast. A defect verified only on nodes does not satisfy the hypothesis unless interpolation, coverage, integration and arithmetic errors are also charged.
\end{proposition}
\begin{proof}
At the terminal date the claim is the first defect bound. At an earlier date subtract the exact paired residual identity from \eqref{eq:pairdefect53}, and use that a probability coupling has operator norm one on bounded functions. This gives an error at most $\eta_t+\beta_ta_{t+1}$. The same norm bound applied to an action coupling proves the last assertion. The argument is valid on the reachable pair domains when every child domain is included in the next domain.
\end{proof}

The result supplies an approximation account independent of how $\psi$ is fitted. For example, a verified expansion $\psi_t(x,y)=\sum_{j=1}^{r_t}u_{tj}(x)v_{tj}(y)$ stores factors instead of a square table. On an $N$-state representation it uses $O(r_tN)$ coefficients rather than $O(N^2)$. With a $q$-node common-innovation coupling, evaluating one pair integral costs $O(qr_t)$ factor products after the factor queries. It does \emph{not} in general factor into a product of marginal expectations; that further reduction requires a product coupling or another proved integration identity. Computing and uniformly validating the defects remains work. A small rank without a verified defect cannot be used in a safety claim.

A local cover offers a second implementation. If an integrand is Lipschitz on a certified region with constant $L$, replacing it by its value at a cover center of radius $h$ costs at most $Lh$. A partition of coupling mass with binwise radii $h_j$ costs at most $\sum_jp_jL_jh_j$. When an acquired actor jumps, the region must be split at the relevant cell boundaries or its entire action range enclosed. A global Lipschitz claim for that discontinuous actor is not available. These two procedures give concrete ways to obtain the $\eta$ and $\tau$ in Proposition~\ref{prop:compressed53}; they are additional implementations covered by the theorem, not unperformed experiments reported as scaling evidence.

\begin{proposition}[On-demand pair-tree work]\label{prop:work53}
Suppose a pass evaluates $K_t$ candidate or action-cover pairs on each of $C_t$ acquired cells at date $t$. Each nonterminal pair integration has at most $q$ children, and its last innovation has an analytic integral. Let $h=T-t$ and let $\kappa_d$ bound the work of a primitive pair node, including actor queries. The arithmetic work of a complete pass is at most a constant times
\begin{equation}\label{eq:work53}
 \sum_{t<T}C_tK_t\kappa_d\sum_{j=0}^{h-1}q^j.
\end{equation}
Streaming child bins depth first uses $O(hbd)$ live numerical coordinates for a batch of $b$ pairs, in addition to the stored policy representation, query index and durable output. It does not store the Cartesian square of the acquired state set.
\end{proposition}
\begin{proof}
There are at most $q^j$ nodes at depth $j$ for each initial pair. Summing depths and initial pairs gives \eqref{eq:work53}; identity pruning can only reduce work. Depth-first evaluation retains the current parent and its accumulated sum at each level, giving the memory bound. Each level stores states, their separately propagated difference, actions and finitely many interval accumulators, all of dimension $O(bd)$.
\end{proof}

For the executed two-dimensional $n\times n$ acquired grid, each date's integer actor is indexed by dyadic rectangle minima and maxima. Preprocessing and storage are $O(n^2\log^2n)$, and a rectangle query uses four entries for each extremum. The exact action quantum makes these extrema integer operations. In $d$ dimensions an analogous tensor index has $O(n^d\log^dn)$ entries and $2^d$ corners per extremum. Accordingly, the present implementation saves pair-state storage but remains tensor in state acquisition and exponential in the uncompressed horizon. Its short-horizon work is measured rather than asserted to be negligible.

For an explicit accuracy account in a smooth, verified region, a state cover with radius $h_x$ uses $O(h_x^{-d})$ cells; an $m$-dimensional action cover uses $O(h_a^{-m})$ elements; and an $s$-dimensional innovation cover uses $O(h_z^{-s})$ bins. Substituting these counts in \eqref{eq:work53} and allocating their certified modulus errors through \eqref{eq:pairerror53} gives a terminating finite-accuracy construction. These worst-case exponents are costs of that construction, not lower bounds for all representations. The original separated state/action/integration resource allocation and its lower bounds remain in the complete development and historical article. The new signed certificate is an additional use of that resource account.

The raw records count primitive pair nodes, coupling bins, analytic terminal pairs, rectangle queries, ambiguous actor boxes, index bytes, peak process memory, and cumulative clocks. The scalar-width and symmetric-contrast gates are additionally evaluated on the same proposal comparisons. Their acceptance counts compare certificates holding the proposals fixed; they are not independently timed alternative full-policy solvers. In particular, this ablation cannot support a claim that the pair recursion costs less than every possible policy-evaluation method.

\section{Approximate nullity and model error}\label{sec:null53}
Exact action-nullity is a useful boundary case, not a premise to impose on a learned error. Suppose an affine action exposure is known,
$F(x,a,z)=f(x,z)+Ba$, and a represented component is $n(y)=\phi(Ly)$, where $\phi$ has a verified Lipschitz constant $K$. Under the same innovation law for both actions,
\begin{equation}\label{eq:projected53}
 |(P^a-P^b)n(x)|\le K\|LB(a-b)\|.
\end{equation}
This follows pointwise and hence after integration. If $LB=0$, cancellation is exact. If not, the displayed positive number is charged. A neural realization can provide $K$ by a product of certified layer norms; a sum of rectified affine units can instead sum the absolute output coefficients times the corresponding directional exposures. These are bounds for the serialized function, not an inference that its unknown fitting error has the same structure. A native representation of that function has the same guarantee.

\begin{proposition}[Perturbation allowances]\label{prop:null53}
Suppose $n=\phi\circ L$ as above, $\|L(B-\widehat B)\|\le\delta_B$, and the innovation contribution is $GZ$. Let the two conditional innovation laws admit a coupling with $\mathbb E\|Z_a-Z_b\|\le\delta_Z$. Then
\begin{equation}\label{eq:misspec53}
 |(P^a-P^b)n(x)|\le
 K\{\|L\widehat B(a-b)\|+\delta_B\|a-b\|+\|LG\|\delta_Z\}.
\end{equation}
A represented remainder with uniform absolute error $\delta_n$ adds at most $2\delta_n$. For $n_M(y)=M(y_1-(2+\eta)y_2+1)_+$ in the unperturbed investment economy,
\begin{equation}\label{eq:eta53}
 |(P^a-P^b)n_M(x)|\le M|\eta|\,|a-b|/4.
\end{equation}
\end{proposition}
\begin{proof}
Couple the innovations as stipulated and subtract their projected successor states. The triangle inequality and the Lipschitz property bound the difference in $n$ by the right side of \eqref{eq:misspec53}; integration preserves the bound. A uniformly bounded remainder contributes at most its range $2\delta_n$. Finally $(1,-2-\eta)(1/2,1/4)'=-\eta/4$, which proves \eqref{eq:eta53}.
\end{proof}

The error component must actually be identified as part of $J^\pi-h$. One legitimate use is to evaluate a corrected critic $h+n$ by the residual construction and then add the verified contrast of $n$. Neither a sampled small projection nor a trained weight vector is itself a certificate of the unknown remainder. Equation~\eqref{eq:misspec53} also distinguishes kernel misspecification from a change in counterfactual correlation: if each action's marginal innovation law is unchanged, the economic expectation is unchanged. The diagnostic below changes conditional means, not merely the coupling used to compare identical marginals.

There is a concrete false-null failure in the original economy. At the final decision date take $x=(0,0)$, incumbent $b=1/8$, proposal $a=0$, $M=4096$ and $\eta=1/16$. Direct rational evaluation gives a strictly positive true advantage. The nuisance contrast is $8$, so the critic $h_T=g-n_M$ reports the true advantage minus $\beta8$, which is negative. A zero error allowance would accept a harmful action. Charging \eqref{eq:eta53} rejects that conclusion. The deposited exact-arithmetic test verifies both signs; the full-cell diagnostic tests the same phenomenon with nonzero exposure and innovation-mean perturbations.
